\element{interface}{
  \begin{questionmult}{interface}\bareme{haut=2,p=0}
    A une interface entre deux milieux, quelles sont les grandeurs qui sont toujours continues ?
      \begin{reponses}
        \bonne{La pression}
        \bonne{La vitesse}
        \mauvaise{La composante du champ électrique normale à l'interface}
        \bonne{La composante du champ électrique tangente à l'interface}
      \end{reponses}
  \end{questionmult}
}
%\setgroupmode{vitesses}{cyclic}

\element{adaptationImpedance}{
  \begin{questionmult}{adaptationImpedance}\bareme{haut=2,p=0}
    Lorsqu'il y a adaptation d'impédance entre deux fluides,
    %\begin{multicols}{2}
      \begin{reponses}
        \bonne{les impédances acoustiques des fluides sont égales.}
        \mauvaise{le coefficient de réflexion est égal à 1}
        \bonne{le coefficient de transmission est égal à 1}
      \end{reponses}
    %\end{multicols}
  \end{questionmult}
}

\element{laplace}{
  \begin{questionmult}{laplace}\bareme{haut=2,p=0}
    Parmi les expressions suivantes de force de Laplace pour une distribution linéique, volumique ou surfacique, lesquelles sont correctes ?
    \begin{multicols}{2}
      \begin{reponses}
        \bonne{$\vv{j_S}\wedge\vv{B}dS$}
        \bonne{$\vv{j}\wedge\vv{B}dV$}
        \bonne{$i\vv{dl}\wedge\vv{B}$}
      \end{reponses}
    \end{multicols}
  \end{questionmult}
}

\element{courantSurfacique}{
  \begin{question}{courantSurfacique}\bareme{1}
    Quelle est l'unité du vecteur densité surfacique de courant électrique ?
    \begin{multicols}{4}
      \begin{reponses}
        \bonne{$\unit{A.m^{-1}}$}
        \mauvaise{$\unit{C.m^{-2}.s^{-1}}$}
        \mauvaise{$\unit{A}$}
        \mauvaise{$C.s^{-1}$}
      \end{reponses}
    \end{multicols}
  \end{question}
}